% ============================================================
%  Simulation Results: Scenarios A and B
%  Thesis Chapter – Results Section
% ============================================================
%
% Compilation note:  This file should be \input{} or \include{}d from
% the main thesis document.  Required packages (assumed loaded in preamble):
%   booktabs, multirow, siunitx, amsmath, hyperref, geometry, graphicx
% ============================================================

\chapter{Simulation Results: Scenarios A and B}
\label{chap:results_AB}

This chapter presents the numerical results of the simulation study for Scenario~A
(Single Supplier Mandate) and Scenario~B (Multiple Suppliers with Independent
Balancing Groups).  Both scenario families are evaluated in a baseline variant
(\emph{without} a Renewable Energy Community, REC) and an enhanced variant
(\emph{with} REC), yielding four primary sub-scenarios: A1, A2, B1, and B2.
An additional pair of \emph{mixed-portfolio} sub-variants (suffixed ``-M'') tests
robustness to alternative customer-to-supplier assignments.

All simulations cover the 2016 calendar year (35\,040 quarter-hourly
intervals, $\Delta t = 15\,\text{min}$) on the SimBench low-voltage rural
network \texttt{1-LV-rural3-{}-2-no\_sw} located near Vienna, Austria
(48.31°\,N, 14.29°\,E).  The community comprises \textbf{nine metered nodes}:
three prosumers (photovoltaic generation + load) and six pure consumers.
Retail electricity is billed at a time-varying tariff; feed-in compensation is
fixed at €82.40\,/\,MWh throughout the year.  Day-ahead (DA), intra-day (ID),
and imbalance prices are sourced from the Austrian APCS/EPEX market data for
2016.

The KPIs used for evaluation are defined in Chapter~\ref{sec:economic_kpis}
and summarised in Table~\ref{tab:kpi_summary}.


% ============================================================
\section{Scenario A -- Single Supplier Mandate}
\label{sec:results_A}
% ============================================================

Scenario~A represents the regulatory configuration in which all community
members -- regardless of whether they are consumers or prosumers -- are contractually
bound to a \textbf{single energy supplier} (Supplier~A, Balancing Group~BG\_A).
This eliminates any competitive market dynamics at the retail level but
simultaneously maximises portfolio aggregation benefits for the supplier's
balancing group.

% ------------------------------------------------------------
\subsection{Scenario A1: Single Supplier, No REC (Net Generators)}
\label{sec:results_A1}

The baseline configuration yields exceptionally stable balancing group performance, with near-perfect alignment between forecast and realised positions throughout the year. Table~\ref{tab:A1_balancing} summarises the annual balancing statistics for Supplier~A's balancing group (BG\_A). The actual metered net position (91.65\,MWh) differs from the intra-day forecast position (91.65\,MWh) by only +0.0075\,MWh, resulting in virtually zero net annual imbalance. The marginal system position is classified as \textsc{short} (deficit), producing a small net balancing settlement of +€0.58 (balancing rewards €71.15 minus penalties €70.57). This near-symmetric outcome reflects the aggregation of demand and generation within the supplier's balancing group, where positive and negative quarter-hourly deviations largely cancel out over the year. Forecasting is relatively straightforward because of the absence of REC sharing. The net generator prosumer mix, with higher installed PV capacity relative to load, yields a lower net energy position (91.65\,MWh) compared to the mixed prosumer configuration, as greater on-site generation offsets more demand.

Figure~\ref{fig:imbalance_annual_results_A1} illustrates the monthly balancing group positions (actual metered versus forecast). Both series track each other closely, exhibiting a pronounced seasonal pattern driven by variations in net load demand. Peak values occur during winter months (November--December), with December recording the annual maximum of approximately 12.50\,MWh (actual 12.5029\,MWh, forecast 12.5087\,MWh). A steady decline follows from January through May, reaching the annual minimum in August (actual 1.37\,MWh, forecast 1.37\,MWh), reflecting strong offset from solar generation in the net generator portfolio. In autumn (September--October), both curves recover, rising from approximately 7.88\,MWh in September to around 8.38\,MWh in October as net load demand increases with shorter days.

\begin{table}[htbp]
\centering
\caption{Scenario A1 (Net Generators) -- Balancing group annual position (Supplier~A / BG\_A).}
\label{tab:A1_balancing}
\renewcommand{\arraystretch}{1.3}
\begin{tabular}{@{}lrl@{}}
\toprule
\textbf{Balancing Metric}          & \textbf{Value}    & \textbf{Unit} \\
\midrule
Actual BG net position             & 91.65             & MWh \\
Forecast BG net position (ID)      & 91.65             & MWh \\
Net imbalance volume               & +0.0075           & MWh \\
System position                    & \multicolumn{2}{l}{\textsc{short (Deficit)}} \\
\addlinespace
Balancing reward (annual)          & 71.15             & € \\
Imbalance penalty (annual)         & 70.57             & € \\
Net balancing settlement           & +0.58             & € \\
\bottomrule
\end{tabular}
\end{table}

\begin{figure}[htbp]
    \centering
    \includegraphics[width=0.8\linewidth]{pdflatex template/img/imbalance-A1_plot.png}
    \caption{Monthly balancing group position and imbalance: Scenario A1 -- Net Generators (Supplier~A).}
    \label{fig:imbalance_annual_results_A1}
\end{figure}

Retail billing dominates the financial outcome. Supplier~A generated annual retail sales revenue of €29{,}814.30 from its customer portfolio, while feed-in compensation to prosumers amounted to €4{,}670.08. Feed-in costs are notably higher than in the mixed prosumer baseline (€3{,}357.64), reflecting the larger net generation capacity of the prosumer portfolio: more surplus energy is exported to the grid, increasing the supplier's feed-in obligation.

Figure~\ref{fig:costsA} decomposes the supplier's total annual costs of €8{,}204.21 into three principal components with contrasting seasonal profiles:

Market purchases (€3{,}464; 42.22\%) represent the largest single expenditure, reflecting wholesale procurement to cover net demand. Costs peak in winter (e.g., December €525, November €494) due to elevated net load demand and minimal solar offset, and reach minima in summer (e.g., August €125, May €151), where high prosumer generation reduces the residual demand that must be procured externally from the day-ahead and intra-day markets. Retail feed-in costs (€4{,}670; 56.92\%) compensate prosumers for surplus exports. These costs exhibit the opposite seasonal pattern, peaking in summer (e.g., May €671, August €612) when photovoltaic surplus is greatest, and approaching lower levels in winter (e.g., December €105, January €134). Notably, in the net generator configuration feed-in costs \emph{exceed} market purchases, unlike the mixed prosumer baseline where market purchases dominated -- this reflects the larger installed PV capacity. Balancing penalties (€70.57; 0.86\%) are negligible, ranging from €3.47 to €8.75 monthly. The slight summer increase reflects marginally higher short-term variability in photovoltaic output. Monthly total costs remain relatively stable (€493--€828), as the opposing seasonal patterns of market purchases and feed-in costs partially offset each other.

\begin{figure}[htbp]
    \centering
    \includegraphics[width=0.8\linewidth]{pdflatex template/img/revenue_componets.png}
    \caption{Monthly revenue components: Scenario A1 -- Net Generators (Supplier~A).}
    \label{fig:revenue_componets}
\end{figure}

\begin{figure}[htbp]
    \centering
    \includegraphics[width=0.8\linewidth]{pdflatex template/img/finacial_overview.png}
    \caption{Monthly revenue, costs, and profit/loss: Scenario A1 -- Net Generators (Supplier~A).}
    \label{fig:finacial_overview1}
\end{figure}

\begin{figure}[htbp]
    \centering
    \includegraphics[width=0.8\linewidth]{pdflatex template/img/cost_A1.png}
    \caption{Monthly cost components: Scenario A1 -- Net Generators (Supplier~A).}
    \label{fig:costsA}
\end{figure}

Figures~\ref{fig:revenue_componets} and \ref{fig:finacial_overview1} illustrate the monthly revenue components and overall profit/loss trajectory, respectively. Revenue is overwhelmingly dominated by retail sales (98.83\%; €29{,}814.30), with energy market sales (€283.12; 0.94\%) and balancing rewards (€71.15; 0.24\%) contributing only marginally. Market sales are higher than in the mixed prosumer configuration (€174.75), consistent with the greater volume of surplus generation available for wholesale export. Costs are primarily driven by retail feed-in payments (56.92\%; €4{,}670.08) and market purchases (42.22\%; €3{,}463.56), with balancing penalties (0.86\%; €70.57) remaining negligible.

The annual profit-and-loss statement (Table~\ref{tab:A1_pnl}) shows total revenue of €30{,}168.57 and total costs of €8{,}204.21, resulting in an annual profit of €21{,}964.36 (monthly average €1{,}830.36; range €1{,}097.03--€2{,}244.91). Overall, the net generator configuration demonstrates excellent forecast alignment, near-zero imbalance exposure, and robust profitability driven overwhelmingly by retail sales, with feed-in payments and wholesale procurement as the primary cost drivers. Compared to the mixed prosumer baseline (profit €29{,}515.86), the net generator scenario yields lower profit (€21{,}964.36) owing to reduced retail sales revenue (fewer consumers/lower load) and higher feed-in costs. Seasonal variations in individual components are largely offset, yielding stable monthly financial performance.

\begin{table}[htbp]
\centering
\caption{Scenario A1 (Net Generators) -- Annual profit and loss statement (Supplier~A).}
\label{tab:A1_pnl}
\renewcommand{\arraystretch}{1.0}
\begin{tabular}{@{}lrr@{}}
\toprule
\textbf{P\&L Component}             & \textbf{Amount (€)} & \textbf{Share (\%)} \\
\midrule
\multicolumn{3}{@{}l}{\textit{Revenue}} \\
\quad Energy market sales            & 283.12              & 0.94 \\
\quad Balancing rewards              & 71.15               & 0.24 \\
\quad Retail sales                   & 29{,}814.30         & 98.83 \\
\textbf{Total Revenue}               & \textbf{30{,}168.57} & \textbf{100.00} \\
\addlinespace
\multicolumn{3}{@{}l}{\textit{Costs}} \\
\quad Energy market purchases        & 3{,}463.56          & 42.22 \\
\quad Imbalance penalties            & 70.57               & 0.86 \\
\quad Retail purchases (feed-in)     & 4{,}670.08          & 56.92 \\
\textbf{Total Costs}                 & \textbf{8{,}204.21}  & \textbf{100.00} \\
\addlinespace
\textbf{Annual Profit}               & \textbf{€21{,}964.36} & \\
Monthly average                      & €1{,}830.36         & \\
Monthly minimum                      & €1{,}097.03         & \\
Monthly maximum                      & €2{,}244.91         & \\
\bottomrule
\end{tabular}
\end{table}

% ------------------------------------------------------------
\subsection{Scenario A1: Single Supplier, No REC (Mixed Prosumers) -- Baseline}
\label{sec:results_A1_mixed}
% ------------------------------------------------------------

The baseline configuration yields exceptionally stable balancing group performance, with near-perfect alignment between forecast and realised positions throughout the year. Table~\ref{tab:A1_mixed_balancing} summarises the annual balancing statistics for Supplier~A's balancing group (BG\_A). The actual metered net position (148.07\,MWh) differs from the intra-day forecast position (148.08\,MWh) by only $-$0.0005\,MWh, resulting in virtually zero net annual imbalance. The marginal system position is classified as \textsc{long} (surplus), producing a small net balancing settlement of +\euro{}0.35 (balancing rewards \euro{}72.08 minus penalties \euro{}71.73). This near-symmetric outcome reflects the aggregation of demand and generation within the supplier's balancing group, where positive and negative quarter-hourly deviations largely cancel out over the year. Forecasting is relatively straightforward because of the absence of REC sharing. Consequently, the combined demand and PV generation profiles partially offset each other, reducing systematic imbalances at the balancing group level.

Figure~\ref{fig:imbalance_annual_results_A1_mixed} illustrates the monthly balancing group positions (actual metered versus forecast). Both series track each other closely, exhibiting a pronounced seasonal pattern driven by variations in net load demand. Peak values occur during winter months (January--February and November--December), with December recording the annual maximum of approximately 18.1\,MWh (actual 18.07\,MWh, forecast 18.07\,MWh). A steady decline follows from March to May, reaching the annual minimum in May (actual 6.57\,MWh, forecast 6.55\,MWh). Summer months (June--August) remain consistently low, with the absolute trough in August (actual 4.97\,MWh, forecast 4.97\,MWh), reflecting strong offset from solar generation. In autumn (September--October), both curves recover, rising from approximately 12.1\,MWh in September to around 12.6\,MWh in October as net load demand increases with shorter days.

\begin{table}[htbp]
\centering
\caption{Baseline Scenario A1 (Mixed Prosumers) -- Balancing group annual position (Supplier~A / BG\_A).}
\label{tab:A1_mixed_balancing}
\renewcommand{\arraystretch}{1.3}
\begin{tabular}{@{}lrl@{}}
\toprule
\textbf{Balancing Metric}          & \textbf{Value}    & \textbf{Unit} \\
\midrule
Actual BG net position             & 148.07            & MWh \\
Forecast BG net position (ID)      & 148.08            & MWh \\
Net imbalance volume               & $-$0.0005         & MWh \\
System position                    & \multicolumn{2}{l}{\textsc{long (Surplus)}} \\
\addlinespace
Balancing reward (annual)          & 72.08             & \euro{} \\
Imbalance penalty (annual)         & 71.73             & \euro{} \\
Net balancing settlement           & +0.35             & \euro{} \\
\bottomrule
\end{tabular}
\end{table}

\begin{figure}[htbp]
    \centering
    \includegraphics[width=0.8\linewidth]{pdflatex template/img/imbalance-A1_mixed_plot.png}
    \caption{Monthly balancing group position and imbalance: Baseline Scenario A1 -- Mixed Prosumers (Supplier~A).}
    \label{fig:imbalance_annual_results_A1_mixed}
\end{figure}

Retail billing dominates the financial outcome. Supplier~A generated annual retail sales revenue of \euro{}37{,}953.39 from six consumers and three prosumers, while feed-in compensation to prosumers amounted to \euro{}3{,}357.64. Feed-in costs remain comparatively modest due to substantial on-site self-consumption, limiting net exports to the grid.

Figure~\ref{fig:costsA_mixed} decomposes the supplier's total annual costs of \euro{}8{,}684.37 into three principal components with contrasting seasonal profiles:

Market purchases (\euro{}5{,}255; 60.5\%) represent the largest expenditure, reflecting wholesale procurement to cover net demand. Costs peak in winter (e.g., December \euro{}776, November \euro{}738) due to elevated net load demand and minimal solar offset, and reach minima in summer (e.g., August \euro{}203, May \euro{}237), where high prosumer generation reduces the residual demand that must be procured externally from the day-ahead and intra-day markets. Retail feed-in costs (\euro{}3{,}358; 38.7\%) compensate prosumers for surplus exports. These costs exhibit the opposite seasonal pattern, peaking in summer (e.g., May \euro{}529, August \euro{}472) when photovoltaic surplus is greatest, and approaching negligible levels in winter (e.g., December \euro{}54, January \euro{}74). Balancing penalties (\euro{}71.73; 0.8\%) are negligible, ranging from \euro{}3.76 to \euro{}8.71 monthly. The slight summer increase reflects marginally higher short-term variability in photovoltaic output. Monthly total costs remain relatively stable (\euro{}569--\euro{}862), as the opposing seasonal patterns of market purchases and feed-in costs partially offset each other.

Figures~\ref{fig:revenue_componets_mixed} and \ref{fig:finacial_overview_mixed} illustrate the monthly revenue components and overall profit/loss trajectory, respectively. Revenue is overwhelmingly dominated by retail sales (99.35\%; \euro{}37{,}953.39), with energy market sales (\euro{}174.75; 0.46\%) and balancing rewards (\euro{}72.08; 0.19\%) contributing only marginally. Costs are primarily driven by market purchases (60.52\%; \euro{}5{,}255.00) and retail feed-in payments (38.66\%; \euro{}3{,}357.64), with balancing penalties (0.83\%; \euro{}71.73) remaining negligible.

\begin{figure}[htbp]
    \centering
    \includegraphics[width=0.8\linewidth]{pdflatex template/img/revenue_componets_mixed.png}
    \caption{Monthly revenue components: Baseline Scenario A1 -- Mixed Prosumers (Supplier~A).}
    \label{fig:revenue_componets_mixed}
\end{figure}

\begin{figure}[htbp]
    \centering
    \includegraphics[width=0.8\linewidth]{pdflatex template/img/finacial_overview_mixed.png}
    \caption{Monthly revenue, costs, and profit/loss: Baseline Scenario A1 -- Mixed Prosumers (Supplier~A).}
    \label{fig:finacial_overview_mixed}
\end{figure}

\begin{figure}[htbp]
    \centering
    \includegraphics[width=0.8\linewidth]{pdflatex template/img/cost_A1_mixed.png}
    \caption{Monthly cost components: Baseline Scenario A1 -- Mixed Prosumers (Supplier~A).}
    \label{fig:costsA_mixed}
\end{figure}

The annual profit-and-loss statement (Table~\ref{tab:A1_mixed_pnl}) shows total revenue of \euro{}38{,}200.22 and total costs of \euro{}8{,}684.37, resulting in an annual profit of \euro{}29{,}515.86 (monthly average \euro{}2{,}459.65; range \euro{}1{,}512.67--\euro{}3{,}160.92). Overall, the baseline configuration demonstrates excellent forecast alignment, near-zero imbalance exposure, and robust profitability driven overwhelmingly by retail sales, with wholesale procurement and feed-in payments as the primary cost drivers. Seasonal variations in individual components are largely offset, yielding stable monthly financial performance.

\begin{table}[htbp]
\centering
\caption{Baseline Scenario A1 (Mixed Prosumers) -- Annual profit and loss statement (Supplier~A).}
\label{tab:A1_mixed_pnl}
\renewcommand{\arraystretch}{1.0}
\begin{tabular}{@{}lrr@{}}
\toprule
\textbf{P\&L Component}             & \textbf{Amount (€)} & \textbf{Share (\%)} \\
\midrule
\multicolumn{3}{@{}l}{\textit{Revenue}} \\
\quad Energy market sales            & 174.75              & 0.46 \\
\quad Balancing rewards              & 72.08               & 0.19 \\
\quad Retail sales                   & 37{,}953.39         & 99.35 \\
\textbf{Total Revenue}               & \textbf{38{,}200.22} & \textbf{100.00} \\
\addlinespace
\multicolumn{3}{@{}l}{\textit{Costs}} \\
\quad Energy market purchases        & 5{,}255.00          & 60.52 \\
\quad Imbalance penalties            & 71.73               & 0.83 \\
\quad Retail purchases (feed-in)     & 3{,}357.64          & 38.66 \\
\textbf{Total Costs}                 & \textbf{8{,}684.37}  & \textbf{100.00} \\
\addlinespace
\textbf{Annual Profit}               & \textbf{\euro{}29{,}515.86} & \\
Monthly average                      & \euro{}2{,}459.65         & \\
Monthly minimum                      & \euro{}1{,}512.67         & \\
Monthly maximum                      & \euro{}3{,}160.92         & \\
\bottomrule
\end{tabular}
\end{table}

% ------------------------------------------------------------
\subsection{Intra-Scenario Comparison: A1 versus A2}
\label{sec:results_A_comparison}
% ------------------------------------------------------------

Table~\ref{tab:A_comparison} presents a side-by-side KPI comparison for the
two A sub-scenarios.

\begin{table}[htbp]
\centering
\caption{Scenario A -- Key Performance Indicator comparison: A1 vs.\ A2.}
\label{tab:A_comparison}
\renewcommand{\arraystretch}{1.35}
\begin{tabular}{@{}lrrr@{}}
\toprule
\textbf{KPI}              & \textbf{A1 (No REC)}   & \textbf{A2 (With REC)}  & \textbf{$\Delta$ (A2$-$A1)} \\
\midrule
Total revenue (€)         & 75{,}794.55            & 62{,}624.97             & $-13{,}169.58$ \\
Total costs (€)           & 7{,}636.82             & 10{,}441.26             & $+2{,}804.44$ \\
Annual profit (€)         & 68{,}157.72            & 52{,}183.72             & $-15{,}974.00$ \\
Profit change (\%)        & ---                    & ---                     & $-23.4$\,\% \\
\addlinespace
Retail sales revenue (€)  & 68{,}879.99            & 62{,}407.61             & $-6{,}472.38$ \\
Balancing reward (€)      & 6{,}853.84             & 210.25                  & $-6{,}643.59$ \\
Imbalance penalty (€)     & 678.08                 & 212.28                  & $-465.80$ \\
Feed-in cost (€)          & 2{,}412.54             & 153.35                  & $-2{,}259.19$ \\
Market purchases (€)      & 4{,}546.20             & 10{,}075.62             & $+5{,}529.42$ \\
\addlinespace
Total imbalance (MWh)     & 177.02                 & $-0.04$                 & $-177.06$ \\
System position           & \textsc{short}         & \textsc{long}           & --- \\
REC shared energy (MWh)   & ---                    & 75.88                   & --- \\
\addlinespace
Monthly avg.\ profit (€)  & 5{,}679.81             & 4{,}348.64              & $-1{,}331.17$ \\
Monthly min.\ profit (€)  & 3{,}402.22             & 2{,}424.90              & $-977.32$ \\
Monthly max.\ profit (€)  & 9{,}489.50             & 6{,}880.52              & $-2{,}608.98$ \\
\bottomrule
\end{tabular}
\end{table}

The KPI table reveals a \textbf{supplier profit penalty associated with
introducing a REC} in the single-supplier context.  The total annual profit
decreases by €15{,}974.00 ($-23.4$\,\%), driven primarily by two mechanisms:
(i) the elimination of the opportunistic balancing reward, and
(ii) the structural reduction in retail grid sales as community self-consumption
partially displaces external procurement.

From the supplier's perspective, the REC is cost-reducing only in the dimension
of feed-in obligations (saving €2{,}259.19).  However, this is more than offset
by the combined loss of balancing income and retail margin.  The result
underscores a fundamental tension in the supply business under energy community
regulation: while the REC benefits consumers through reduced grid fees and the
community through lower net import, the supplier bears the financial consequence
of reduced throughput.

Seasonally, the profit reduction is most severe in summer months (highest PV
generation, maximum shared energy), where the monthly maximum drops from
€9{,}489.50 to €6{,}880.52 ($-27.5$\,\%).  In winter, where PV output is low
and shared energy minimal, the REC has little effect.

% ============================================================
\section{Scenario B -- Multiple Suppliers with Independent Balancing Groups}
\label{sec:results_B}
% ============================================================

Scenario~B relaxes the single-supplier mandate and introduces \textbf{two
independent energy suppliers}, each managing its own balancing group under
separate forecasting and settlement responsibilities.  The customer allocation
is:
\begin{itemize}
    \item \textbf{Supplier~A / BG\_SUP\_A\_001}: six pure consumers (no on-site
          generation).
    \item \textbf{Supplier~B / BG\_SUP\_B\_001}: three prosumers (PV + load).
\end{itemize}

This segregation reflects a realistic market scenario in which consumers and
prosumers independently choose their preferred supplier.  Supplier~A operates a
\emph{pure load} portfolio, while Supplier~B operates a \emph{mixed
prosumer} portfolio with intermittent net generation.

% ------------------------------------------------------------
\subsection{Scenario B1: Baseline -- Multiple Suppliers Without REC}
\label{sec:results_B1}
% ------------------------------------------------------------

\subsubsection{Supplier A (Consumer Portfolio)}

Supplier~A's portfolio is entirely demand-side, with no on-site generation.
The balancing group therefore maintains a consistently positive net load
position; energy market sales are zero for the year.

\begin{table}[htbp]
\centering
\caption{Scenario B1 -- Annual P\&L: Supplier~A (consumer-only portfolio).}
\label{tab:B1_supA_pnl}
\renewcommand{\arraystretch}{1.3}
\begin{tabular}{@{}lrr@{}}
\toprule
\textbf{P\&L Component}             & \textbf{Amount (€)} & \textbf{Share (\%)} \\
\midrule
\multicolumn{3}{@{}l}{\textit{Revenue}} \\
\quad Energy market sales            & 0.00                & 0.00 \\
\quad Balancing rewards              & 176.78              & 0.35 \\
\quad Retail sales                   & 49{,}958.69         & 99.65 \\
\textbf{Total Revenue}               & \textbf{50{,}135.47} & \textbf{100.00} \\
\addlinespace
\multicolumn{3}{@{}l}{\textit{Costs}} \\
\quad Energy market purchases        & 8{,}128.50          & 97.86 \\
\quad Imbalance penalties           & 177.18              & 2.13  \\
\quad Retail purchases (feed-in)     & 0.00                & 0.00  \\
\textbf{Total Costs}                 & \textbf{8{,}305.68}  & \textbf{100.00} \\
\addlinespace
\textbf{Annual Profit / Loss}        & \textbf{€41{,}829.79} & \\
Monthly average                      & €3{,}485.82        & \\
Monthly minimum                      & €2{,}153.64        & \\
Monthly maximum                      & €5{,}556.02        & \\
\bottomrule
\end{tabular}
\end{table}

Supplier~A's balancing performance is exceptionally clean:
actual position 248.55\,MWh, forecast 248.50\,MWh, net imbalance +0.05\,MWh
(system: SHORT, Deficit) -- an almost perfect year-long forecast.
Consequently, balancing rewards and penalties are nearly symmetric
(€176.78 vs.\ €177.18), yielding a net balancing loss of only $-$€0.40.
The absence of generation assets means zero energy market sales and zero
feed-in costs.  Retail sales at €49{,}958.69 are the sole revenue source.

\subsubsection{Supplier B (Prosumer Portfolio)}

Supplier~B's three-prosumer portfolio exhibits substantially higher volatility
due to weather-dependent PV generation.

\begin{table}[htbp]
\centering
\caption{Scenario B1 -- Annual P\&L: Supplier~B (prosumer portfolio).}
\label{tab:B1_supB_pnl}
\renewcommand{\arraystretch}{1.3}
\begin{tabular}{@{}lrr@{}}
\toprule
\textbf{P\&L Component}             & \textbf{Amount (€)} & \textbf{Share (\%)} \\
\midrule
\multicolumn{3}{@{}l}{\textit{Revenue}} \\
\quad Energy market sales            & 252.41              & 1.22 \\
\quad Balancing rewards              & 112.73              & 0.54 \\
\quad Retail sales                   & 20{,}364.02         & 98.24 \\
\textbf{Total Revenue}               & \textbf{20{,}729.16} & \textbf{100.00} \\
\addlinespace
\multicolumn{3}{@{}l}{\textit{Costs}} \\
\quad Energy market purchases        & 2{,}604.14          & 50.83 \\
\quad Imbalance penalties           & 107.26              & 2.09  \\
\quad Retail purchases (feed-in)     & 2{,}412.54          & 47.10 \\
\textbf{Total Costs}                 & \textbf{5{,}123.93}  & \textbf{100.00} \\
\addlinespace
\textbf{Annual Profit / Loss}        & \textbf{€15{,}605.23} & \\
Monthly average                      & €1{,}300.44        & \\
Monthly minimum                      & €692.65            & \\
Monthly maximum                      & €1{,}794.55        & \\
\bottomrule
\end{tabular}
\end{table}

Supplier~B generates a positive but modest annual profit of €15{,}605.23.
The prosumer portfolio introduces DA/ID market sales (€252.41) when net
PV generation exceeds aggregate load.  Feed-in costs (€2{,}412.54)
account for 47.10\,\% of total costs -- the second-largest cost item after
market purchases -- reflecting the obligation to compensate prosumer
grid exports at the fixed tariff.  Balancing is well-controlled:
imbalance = +0.10\,MWh across the year (system: SHORT).

\subsubsection{Combined B1 Performance}

The aggregate economic performance of both suppliers is:

\begin{table}[htbp]
\centering
\caption{Scenario B1 -- Combined annual performance of Supplier~A and
         Supplier~B.}
\label{tab:B1_combined}
\renewcommand{\arraystretch}{1.3}
\begin{tabular}{@{}lrrr@{}}
\toprule
\textbf{KPI}           & \textbf{Supplier A}  & \textbf{Supplier B}  & \textbf{Combined} \\
\midrule
Total revenue (€)      & 50{,}135.47          & 20{,}729.16          & 70{,}864.63 \\
Total costs (€)        & 8{,}305.68           & 5{,}123.93           & 13{,}429.61 \\
Annual profit (€)      & 41{,}829.79          & 15{,}605.23          & 57{,}435.02 \\
Market purchases (€)   & 8{,}128.50           & 2{,}604.14           & 10{,}732.64 \\
Feed-in costs (€)      & 0.00                 & 2{,}412.54           & 2{,}412.54  \\
Imbalance penalty (€)  & 177.18               & 107.26               & 284.44      \\
Balancing reward (€)   & 176.78               & 112.73               & 289.51      \\
Total imbalance (MWh)  & 0.05                 & 0.10                 & 0.15        \\
\bottomrule
\end{tabular}
\end{table}

At a combined level, the B1 market structure generates a total profit of
\textbf{€57{,}435.02} across both suppliers, with Supplier~A contributing
72.8\,\% of combined profit despite being the pure-consumer supplier.  The
prosumer supplier (B) is less profitable in absolute terms due to the
feed-in obligation that erodes margin.  Both balancing groups display
remarkably precise forecast-actual alignment ($\leq 0.10$\,MWh annual
imbalance), in stark contrast to A1 where systematic under-forecasting was
evident.

% ------------------------------------------------------------
\subsection{Scenario B2: Multiple Suppliers With REC}
\label{sec:results_B2}
% ------------------------------------------------------------

In B2, two independent RECs are introduced -- one for each balancing group --
enabling internal energy sharing prior to supplier billing.  The REC incentives
are identical to A2 (shared energy bonus: €10\,/\,MWh; self-consumption:
€15\,/\,MWh).  The settlement order is:
DA $\to$ ID $\to$ REC settlement $\to$ Balancing $\to$ Billing.

\subsubsection{REC Energy Sharing Performance}

Both Supplier~A's and Supplier~B's RECs record a total shared energy of
\textbf{37.20\,MWh} per balancing group over the year.  This value represents
the intra-community PV generation absorbed by community loads within each
15-minute interval without transiting the distribution grid.  The combined
community-wide shared energy is therefore 74.40\,MWh -- comparable to A2's
75.88\,MWh, confirming consistency of the physical sharing mechanism across
scenarios.

\subsubsection{Supplier A (Consumer Portfolio with REC)}

\begin{table}[htbp]
\centering
\caption{Scenario B2 -- Annual P\&L: Supplier~A (consumer portfolio, REC).}
\label{tab:B2_supA_pnl}
\renewcommand{\arraystretch}{1.3}
\begin{tabular}{@{}lrr@{}}
\toprule
\textbf{P\&L Component}             & \textbf{Amount (€)} & \textbf{Share (\%)} \\
\midrule
\multicolumn{3}{@{}l}{\textit{Revenue}} \\
\quad Energy market sales            & 0.00                & 0.00 \\
\quad Balancing rewards              & 690.34              & 1.60 \\
\quad Retail sales                   & 42{,}481.18         & 98.40 \\
\textbf{Total Revenue}               & \textbf{43{,}171.52} & \textbf{100.00} \\
\addlinespace
\multicolumn{3}{@{}l}{\textit{Costs}} \\
\quad Energy market purchases        & 7{,}487.67          & 87.73 \\
\quad Imbalance penalties           & 1{,}046.17          & 12.26 \\
\quad Retail purchases (feed-in)     & 0.00                & 0.00  \\
\textbf{Total Costs}                 & \textbf{8{,}533.84}  & \textbf{100.00} \\
\addlinespace
\textbf{Annual Profit / Loss}        & \textbf{€34{,}637.68} & \\
Monthly average                      & €2{,}886.47        & \\
Monthly minimum                      & €1{,}154.97        & \\
Monthly maximum                      & €5{,}658.49        & \\
\bottomrule
\end{tabular}
\end{table}

Supplier~A's retail sales fall by €7{,}477.51 ($-14.9$\,\%) from B1 to B2
due to REC-mediated reduction in consumers' net grid imports.  The REC
introduces a modest but non-trivial imbalance penalty (€1{,}046.17 vs.\ €177.18
in B1) alongside a higher balancing reward (€690.34 vs.\ €176.78), resulting
in a net balancing loss of $-$€355.83.  This deterioration in balancing
performance is attributed to the interaction between REC corrections and the
separate (consumer-only) balancing group's forecast, which does not directly
model the partial offset from prosumer generation.

The system position is \textsc{long (Surplus)} for BG~A in B2:
actual position 211.35\,MWh vs.\ forecast 228.93\,MWh, imbalance $-17.58$\,MWh.
This reversal (from SHORT in B1) occurs because the REC settlement reduces
the consumer group's net import in the actual data but not in the DA/ID
forecast, creating a systematic over-forecast of demand.

\subsubsection{Supplier B (Prosumer Portfolio with REC)}

\begin{table}[htbp]
\centering
\caption{Scenario B2 -- Annual P\&L: Supplier~B (prosumer portfolio, REC).}
\label{tab:B2_supB_pnl}
\renewcommand{\arraystretch}{1.3}
\begin{tabular}{@{}lrr@{}}
\toprule
\textbf{P\&L Component}             & \textbf{Amount (€)} & \textbf{Share (\%)} \\
\midrule
\multicolumn{3}{@{}l}{\textit{Revenue}} \\
\quad Energy market sales            & 0.00                & 0.00 \\
\quad Balancing rewards              & 4{,}193.36          & 15.74 \\
\quad Retail sales                   & 22{,}439.14         & 84.26 \\
\textbf{Total Revenue}               & \textbf{26{,}632.50} & \textbf{100.00} \\
\addlinespace
\multicolumn{3}{@{}l}{\textit{Costs}} \\
\quad Energy market purchases        & 0.00                & 0.00 \\
\quad Imbalance penalties           & 393.21              & 66.53 \\
\quad Retail purchases (feed-in)     & 197.83              & 33.49 \\
\textbf{Total Costs}                 & \textbf{591.03}     & \textbf{100.00} \\
\addlinespace
\textbf{Annual Profit / Loss}        & \textbf{€26{,}041.47} & \\
Monthly average                      & €2{,}170.12        & \\
Monthly minimum                      & €1{,}346.74        & \\
Monthly maximum                      & €3{,}013.91        & \\
\bottomrule
\end{tabular}
\end{table}

Supplier~B's performance is qualitatively transformed by REC introduction.
Wholesale market purchases \emph{drop to zero} (from €2{,}604.14 in B1):
after REC settlement, the corrected prosumer net position is consistently
long, eliminating any requirement to buy from the wholesale market.
Feed-in costs also fall to €197.83 ($-91.8$\,\%) as shared energy absorbs
most of the prosumers' generation surplus.

However, the balancing market position reveals an important structural
imbalance: actual BG position = 109.24\,MWh (net net~load), but the ID
forecast = 0.00\,MWh, yielding an annual imbalance of +109.24\,MWh.
This large discrepancy arises because B2's forecast for the prosumer BG is
constructed on corrected meter readings that approach zero, but the actual
residual consumption is non-trivial in hours with low PV output.  Consequently,
Supplier~B earns an exceptionally large balancing reward of \textbf{€4{,}193.36}
-- the largest balancing income of any individual supplier in the study --
as the short system repeatedly rewards the supplier's inadvertent
\emph{long} imbalance contribution.  This is structurally analogous to A1's
balancing windfall, but operating via a different mechanism.

\subsubsection{Combined B2 Performance}

\begin{table}[htbp]
\centering
\caption{Scenario B2 -- Combined annual performance of Supplier~A and
         Supplier~B.}
\label{tab:B2_combined}
\renewcommand{\arraystretch}{1.3}
\begin{tabular}{@{}lrrr@{}}
\toprule
\textbf{KPI}           & \textbf{Supplier A}  & \textbf{Supplier B}  & \textbf{Combined} \\
\midrule
Total revenue (€)      & 43{,}171.52          & 26{,}632.50          & 69{,}804.02 \\
Total costs (€)        & 8{,}533.84           & 591.03               & 9{,}124.87  \\
Annual profit (€)      & 34{,}637.68          & 26{,}041.47          & 60{,}679.15 \\
Market purchases (€)   & 7{,}487.67           & 0.00                 & 7{,}487.67  \\
Feed-in costs (€)      & 0.00                 & 197.83               & 197.83      \\
Imbalance penalty (€)  & 1{,}046.17           & 393.21               & 1{,}439.38  \\
Balancing reward (€)   & 690.34               & 4{,}193.36           & 4{,}883.70  \\
Total imbalance (MWh)  & $-17.58$             & 109.24               & 91.66       \\
REC shared energy (MWh)& 37.20                & 37.20                & 74.40       \\
\bottomrule
\end{tabular}
\end{table}

The combined B2 profit of \textbf{€60{,}679.15} exceeds B1's €57{,}435.02 by
\textbf{€3{,}244.13 (+5.6\,\%)}.  Unlike the A scenario where REC introduction
decreased supplier profit, the B scenario shows a modest \emph{gain} from
REC introduction.  This positive effect is driven by the remarkable reduction
in feed-in costs (from €2,412.54 to €197.83, saving €2,214.71) and the
large balancing windfall earned by Supplier~B (€4,193.36 vs.\ €112.73 in B1).
Combined imbalance penalties also increase (from €284.44 to €1,439.38), but
the net effect remains positive.

% ------------------------------------------------------------
\subsection{Intra-Scenario Comparison: B1 versus B2}
\label{sec:results_B_comparison}
% ------------------------------------------------------------

\begin{table}[htbp]
\centering
\caption{Scenario B -- KPI comparison B1 vs.\ B2 (combined totals and
         per-supplier breakdown).}
\label{tab:B_comparison}
\renewcommand{\arraystretch}{1.35}
\begin{tabular}{@{}lrrrrrr@{}}
\toprule
 & \multicolumn{2}{c}{\textbf{Supplier A}} & \multicolumn{2}{c}{\textbf{Supplier B}} & \multicolumn{2}{c}{\textbf{Combined}} \\
\cmidrule(lr){2-3}\cmidrule(lr){4-5}\cmidrule(lr){6-7}
\textbf{KPI (€)}        & \textbf{B1}   & \textbf{B2}   & \textbf{B1}   & \textbf{B2}   & \textbf{B1}   & \textbf{B2} \\
\midrule
Annual profit           & 41{,}830      & 34{,}638      & 15{,}605      & 26{,}041      & 57{,}435      & 60{,}679 \\
Total revenue           & 50{,}135      & 43{,}172      & 20{,}729      & 26{,}633      & 70{,}865      & 69{,}804 \\
Total costs             & 8{,}306       & 8{,}534       & 5{,}124       & 591           & 13{,}430      & 9{,}125  \\
Retail sales            & 49{,}959      & 42{,}481      & 20{,}364      & 22{,}439      & 70{,}323      & 64{,}920 \\
Market purchases        & 8{,}129       & 7{,}488       & 2{,}604       & 0             & 10{,}733      & 7{,}488  \\
Feed-in costs           & 0             & 0             & 2{,}413       & 198           & 2{,}413       & 198      \\
Balancing rewards       & 177           & 690           & 113           & 4{,}193       & 290           & 4{,}884  \\
Imbalance penalties     & 177           & 1{,}046       & 107           & 393           & 284           & 1{,}439  \\
\addlinespace
Imbalance (MWh)         & 0.05          & $-17.58$      & 0.10          & 109.24        & 0.15          & 91.66   \\
Shared energy (MWh)     & ---           & 37.20         & ---           & 37.20         & ---           & 74.40   \\
\bottomrule
\end{tabular}
\end{table}

For \textbf{Supplier A}, REC introduction reduces annual profit from
€41{,}829.79 to €34{,}637.68 ($-$€7{,}192.11, $-17.2$\,\%).  This is the
consumer-only supplier: the REC reduces the quantity of electricity it sells
at the retail tariff, directly eroding its retail margin with no compensating
revenue stream.

For \textbf{Supplier B}, REC introduction dramatically \emph{increases} annual
profit from €15{,}605.23 to €26{,}041.47 (+€10{,}436.24, +\textbf{+66.9\,\%}).
The prosumer-only supplier benefits from (i) near-total elimination of
feed-in costs ($-$€2{,}214.71), (ii) elimination of wholesale purchases
($-$€2{,}604.14 saving), and (iii) a large balancing windfall of €4{,}193.36
due to systematic REC-induced over-forecasting.  The net effect is that
Supplier~B's profit share in B2 rises from 27.2\,\% (B1) to 42.9\,\% (B2).

\textbf{Combined}, REC introduction in Scenario~B yields a net gain of +€3{,}244
(+5.6\,\%) but with a marked redistribution of profits: the consumer supplier
loses while the prosumer supplier gains substantially.


% ============================================================
\section{Cross-Scenario Comparison: Scenario A versus Scenario B}
\label{sec:results_AB_crosscomparison}
% ============================================================

This section compares Scenarios~A and~B across all four primary sub-scenarios,
focusing on the system-level aggregate economics, balancing efficiency, and
the incremental impact of REC introduction.

\subsection{Aggregate Supplier Profit Comparison}

Table~\ref{tab:AB_crosscomparison} presents the headline profit figures for all
four configurations side by side.

\begin{table}[htbp]
\centering
\caption{Cross-scenario aggregate profit comparison: A1, A2, B1, B2.}
\label{tab:AB_crosscomparison}
\renewcommand{\arraystretch}{1.35}
\begin{tabular}{@{}lrrrr@{}}
\toprule
\textbf{Metric}                   & \textbf{A1}
                                  & \textbf{A2}
                                  & \textbf{B1}
                                  & \textbf{B2} \\
\midrule
Annual profit -- all suppliers (€)& 68{,}157.72  & 52{,}183.72  & 57{,}435.02  & 60{,}679.15 \\
$\Delta$ vs.\ no-REC baseline (€) & ---           & $-15{,}974$   & ---          & $+3{,}244$ \\
$\Delta$ vs.\ no-REC baseline (\%)& ---           & $-23.4$\,\%   & ---          & $+5.6$\,\% \\
$\Delta$ vs.\ A-scenario (€)      & ---           & ---           & $-10{,}723$  & $+8{,}495$ \\
$\Delta$ vs.\ A-scenario (\%)     & ---           & ---           & $-15.7$\,\%  & $+16.3$\,\% \\
\addlinespace
Total retail revenue (€)         & 68{,}880      & 62{,}408      & 70{,}323     & 64{,}920   \\
Total market purchases (€)       & 4{,}546       & 10{,}076      & 10{,}733     & 7{,}488    \\
Total feed-in costs (€)          & 2{,}413       & 153           & 2{,}413      & 198        \\
Net balancing (€)                & $+6{,}176$    & $-2$          & $+5$         & $+3{,}444$ \\
Shared energy (MWh)              & 0             & 75.88         & 0            & 74.40      \\
\bottomrule
\end{tabular}
\end{table}

\paragraph{Single vs.\ Multiple Suppliers (No REC: A1 vs.\ B1).}
Under the no-REC configuration, the single-supplier mandate (A1) outperforms the
multiple-supplier arrangement (B1) by \textbf{€10{,}723} (+18.7\,\%).  The
advantage is sourced entirely from the anomalously large balancing windfall in
A1 (€6{,}176 net) relative to B1's near-zero net balancing ($+$€5).  In terms
of retail economics -- the structural component of profit -- the two are
comparable (A1: €66{,}467 net retail vs.\ B1: €67{,}910), confirming that both
configurations access similar retail revenues from the same nine customer sites.

The additional €3{,}200 in B1 net retail is attributable to Supplier~B's small
energy market sales (€252.41 from DA/ID net generation), which are not available
to the consolidated A1 portfolio because a net-generator BG offset the consumer
BG's purchase commitment.

\paragraph{REC Impact Across Scenarios (A1$\to$A2 and B1$\to$B2).}
The economic impact of introducing a REC differs markedly between A and B:
\begin{itemize}
    \item In \textbf{Scenario A}, REC introduction reduces total supplier profit
          by €15{,}974 ($-23.4$\,\%).  The single integrated supplier loses
          retail margin on all nine customers simultaneously, while simultaneously
          losing the large (if opportunistic) balancing reward.

    \item In \textbf{Scenario B}, REC introduction increases total supplier profit
          by €3{,}244 (+5.6\,\%).  The benefit is concentrated in Supplier~B's
          prosumer portfolio, where feed-in obligations and wholesale purchases
          are dramatically reduced.
\end{itemize}

This asymmetry has a clear structural explanation.  With a \emph{single} supplier
(A-scenario), the REC reduces total throughput through the supplier's meter --
both the generation and consumption sides of the REC are the same entity's
customers, so the netting reduces billable volume for the sole supplier.  With
\emph{multiple} suppliers (B-scenario), however, the REC crosses supplier
boundaries: shared energy from Supplier~B's prosumers is consumed by
Supplier~A's consumers (or vice versa), reducing each supplier's exposure to
feed-in or procurement costs without necessarily eliminating retail revenue.

\paragraph{Single vs.\ Multiple Suppliers (With REC: A2 vs.\ B2).}

With REC, B2 outperforms A2 by \textbf{€8{,}495} (+16.3\,\%).  The higher
combined profit in B2 is partly due to the asymmetric cross-supplier benefit of
the REC, and partly due to the large balancing windfall earned by Supplier~B
(€4{,}193 vs.\ A2's €210).  However, in A2 total costs are structurally lower
(€10{,}441 vs.\ B2's €9{,}125), reflecting the A-scenario's integrated portfolio
aggregation.

\subsection{Balancing Efficiency Comparison}

\begin{table}[htbp]
\centering
\caption{Balancing efficiency KPIs across all four scenarios.}
\label{tab:balancing_comparison}
\renewcommand{\arraystretch}{1.35}
\begin{tabular}{@{}lrrrr@{}}
\toprule
\textbf{KPI}                     & \textbf{A1}   & \textbf{A2}   & \textbf{B1}   & \textbf{B2 (A+B)} \\
\midrule
Total imbalance (MWh)            & 177.02        & $-0.04$       & 0.15          & 91.66  \\
Total balancing reward (€)       & 6{,}853.84    & 210.25        & 289.51        & 4{,}883.70 \\
Total imbalance penalty (€)      & 678.08        & 212.28        & 284.44        & 1{,}439.38 \\
Net balancing settlement (€)     & $+6{,}175.76$ & $-2.03$       & $+5.07$       & $+3{,}444.32$ \\
Imbalance-to-actual ratio        & 56.5\,\%      & $<0.1$\,\%    & $<0.1$\,\%    & 28.5\,\% \\
System position                  & \textsc{short}& \textsc{long} & \textsc{short}& \textsc{mixed} \\
\bottomrule
\end{tabular}
\end{table}

The balancing results reveal two fundamentally different regimes:

\begin{enumerate}
    \item \textbf{A1 -- high imbalance, high reward:}  A 177.02\,MWh annual
          imbalance (56.5\,\% of actual position) generates €6{,}854 in rewards.
          This is the direct consequence of not correcting the prosumer/consumer
          offset at the BG level: the single BG carries the full prosumer
          generation independently of consumer load, leading to systematic under-
          forecasting of the net position.

    \item \textbf{A2 -- near-perfect balance:}  REC settlement at the 15-minute
          level produces an essentially zero annual imbalance, reducing both
          rewards and penalties to €210 -- a 97\,\% reduction in balancing
          activity.

    \item \textbf{B1 -- excellent balance, low stakes:}  Separating consumers
          (Supplier A) and prosumers (Supplier B) into distinct BGs results in
          very clean forecasts for each group (annual imbalance 0.05 and 0.10\,MWh
          respectively).  Balancing is near-neutral with minimal financial impact.

    \item \textbf{B2 -- imbalance returns through the REC:}  Introducing RECs
          into separate supplier portfolios creates systematic forecast deviations
          (combined imbalance: 91.66\,MWh), because the REC corrections are
          applied to actuals but not fully propagated to ID forecasts.  Supplier~B
          benefits disproportionately (€4{,}193 reward), while Supplier~A incurs
          the largest individual penalty (€1{,}046).
\end{enumerate}

\subsection{REC Shared Energy Impact}

Both A2 and B2 achieve comparable annual shared energy volumes -- A2:
75.88\,MWh, B2: 74.40\,MWh (combined across both BGs).  This confirms that the
physical sharing mechanism is essentially scenario-structure-agnostic: the
available solar generation and community load profiles determine the ceiling of
shareable energy, irrespective of whether that generation and load reside in
the same or different balancing groups.

The economic value extracted from the shared energy, however, differs between
scenarios.  In A2, the 75.88\,MWh saved the supplier €2{,}259\,in feed-in
costs but cost approximately €6,472 in lost retail sales, yielding a net
economic effect of $-$€4,213 from the sharing alone (before balancing effects).
In B2 (combined), the 74.40\,MWh saved €2,215 in feed-in costs (Supplier~B)
and reduced market purchases by €3{,}245 (Supplier~B), while reducing
Supplier~A's retail sales by €7{,}478.  The combined sharing economics are:
$+$€5{,}460 in cost savings vs.\ $-$€5{,}403 in retail revenue reduction --
approximately break-even at the sharing-mechanism level alone, before balancing
effects are included.

\subsection{Mixed-Portfolio Sensitivity Analysis}

To assess robustness of the main scenario results to the assignment of customers
to suppliers, ``mixed-portfolio'' variants were simulated (suffix -M), in which
the prosumer/consumer allocation across balancing groups differs from the primary
configuration.  The key results are:

\begin{table}[htbp]
\centering
\caption{Mixed-portfolio variants -- headline annual profit and balancing KPIs.}
\label{tab:mixed_variants}
\renewcommand{\arraystretch}{1.35}
\begin{tabular}{@{}llrrrl@{}}
\toprule
\textbf{Variant}  & \textbf{Supplier}  & \textbf{Ann.\ Profit (€)}
                  & \textbf{Imbalance (MWh)} & \textbf{Shared (MWh)}
                  & \textbf{System Pos.} \\
\midrule
A1-M              & A (single)          & 13{,}122.47 & $-0.05$  & ---   & \textsc{long}  \\
A2-M              & A (single)          & 11{,}276.47 & $-0.05$  & 90.56 & \textsc{long}  \\
B1-M Supplier A   & A (mixed)           & 5{,}784.51  & 0.00     & ---   & \textsc{short} \\
B1-M Supplier B   & B (mixed)           & 4{,}221.60  & 0.03     & ---   & \textsc{short} \\
B2-M Supplier A   & A (mixed, REC)      & 14{,}812.29 & $-45.28$ & 45.28 & \textsc{long}  \\
B2-M Supplier B   & B (mixed, REC)      & $-3{,}535.83$ & 45.23  & 45.28 & \textsc{short} \\
\bottomrule
\end{tabular}
\end{table}

Several observations emerge:

\begin{enumerate}
    \item The mixed A1/A2 variants are substantially less profitable
          (€13{,}122 and €11{,}276 vs.\ €68{,}158 and €52{,}184) because the
          mixed network contains fewer or smaller customers by design, reducing
          absolute energy throughput.

    \item In the mixed B2 configuration, Supplier~B posts a \textbf{net loss of
          $-$€3{,}535.83}, demonstrating the potential for the cross-supplier REC
          structure to produce financially adverse outcomes for the prosumer
          supplier when the imbalance dynamic works against it.  This highlights
          the supply-risk asymmetry embedded in multi-supplier REC structures.

    \item Shared energy under A2-M reaches 90.56\,MWh -- higher than the
          standard A2's 75.88\,MWh -- indicating that the mixed customer
          composition in the alternative network has a more generation-rich
          profile relative to demand, enabling greater intra-community matching.
\end{enumerate}

The mixed variants confirm that the direction of results is robust across
customer compositions, but the magnitude of financial effects is sensitive to
portfolio mix, particularly in the presence of REC structures.


% ============================================================
\section{Summary of Key Findings}
\label{sec:results_summary}
% ============================================================

Table~\ref{tab:summary_findings} distils the key economic performance indicators
for all four primary sub-scenarios into a single comparative dashboard.

\begin{table}[htbp]
\centering
\caption{Comprehensive KPI dashboard: all four primary scenarios.}
\label{tab:summary_findings}
\renewcommand{\arraystretch}{1.45}
\begin{tabular}{@{}lrrrr@{}}
\toprule
\textbf{KPI}                            & \textbf{A1}       & \textbf{A2}       & \textbf{B1}       & \textbf{B2}       \\
\midrule
\multicolumn{5}{@{}l}{\textit{Scenario structure}} \\
\quad Suppliers                         & 1                 & 1                 & 2                 & 2                 \\
\quad RECs                              & 0                 & 1                 & 0                 & 2                 \\
\quad Prosumers                         & 3                 & 3                 & 3 (Sup.\ B)       & 3 (Sup.\ B)       \\
\quad Consumers                         & 6                 & 6                 & 6 (Sup.\ A)       & 6 (Sup.\ A)       \\
\addlinespace
\multicolumn{5}{@{}l}{\textit{Financial performance (aggregate, €)}} \\
\quad Total revenue                     & 75{,}795          & 62{,}625          & 70{,}865          & 69{,}804          \\
\quad Total costs                       & 7{,}637           & 10{,}441          & 13{,}430          & 9{,}125           \\
\quad \textbf{Annual profit}            & \textbf{68{,}158} & \textbf{52{,}184} & \textbf{57{,}435} & \textbf{60{,}679} \\
\quad Monthly avg.\ profit              & 5{,}680           & 4{,}349           & 4{,}786           & 5{,}057           \\
\addlinespace
\multicolumn{5}{@{}l}{\textit{Revenue decomposition (€)}} \\
\quad Retail sales                      & 68{,}880          & 62{,}408          & 70{,}323          & 64{,}920          \\
\quad Balancing rewards                 & 6{,}854           & 210               & 290               & 4{,}884           \\
\quad Market sales                      & 61                & 7                 & 252               & 0                 \\
\addlinespace
\multicolumn{5}{@{}l}{\textit{Cost decomposition (€)}} \\
\quad Market purchases                  & 4{,}546           & 10{,}076          & 10{,}733          & 7{,}488           \\
\quad Feed-in costs                     & 2{,}413           & 153               & 2{,}413           & 198               \\
\quad Imbalance penalties               & 678               & 212               & 284               & 1{,}439           \\
\addlinespace
\multicolumn{5}{@{}l}{\textit{Balancing KPIs}} \\
\quad Total imbalance (MWh)             & 177.02            & $-0.04$           & 0.15              & 91.66             \\
\quad Net balancing settlement (€)      & $+6{,}176$        & $-2$              & $+5$              & $+3{,}444$        \\
\quad System position                   & Short             & Long              & Short             & Mixed             \\
\addlinespace
\multicolumn{5}{@{}l}{\textit{REC KPIs}} \\
\quad Total shared energy (MWh)         & ---               & 75.88             & ---               & 74.40             \\
\quad Feed-in cost saving (€)           & ---               & $+2{,}259$        & ---               & $+2{,}215$        \\
\quad Retail revenue impact (€)         & ---               & $-6{,}472$        & ---               & $-5{,}403$        \\
\bottomrule
\end{tabular}
\end{table}

The following principal findings are drawn from the simulation:

\begin{enumerate}
    \item \textbf{Single-supplier mandate maximises profit in the no-REC
          baseline (A1):}  The consolidated BG captures the full retail margin
          from all nine customers.  In 2016, it also benefited from an
          asymmetric balancing windfall (€6{,}176 net), which alone accounts for
          the majority of A1's advantage over B1.

    \item \textbf{REC introduction is financially adverse for a single
          integrated supplier (A2):}  The $-23.4$\,\% profit reduction in A2
          stems from (i) reduced retail throughput ($-$€6{,}472) as shared energy
          displaces grid imports, (ii) elimination of the balancing reward, and
          (iii) higher wholesale purchase costs.  The supplier's fee-in saving
          (€2{,}259) only partially offsets these losses.

    \item \textbf{Market fragmentation (B1) reduces aggregate profit but
          improves balancing quality:}  Splitting consumers and prosumers into
          separate BGs (B1) reduces aggregate profit by €10{,}723 versus A1.
          The deficit is structural: the consumer BG cannot net against prosumer
          generation at the DA/ID stage, leading to higher gross market purchases.
          However, both BGs achieve near-perfect forecast alignment.

    \item \textbf{REC introduction under multiple suppliers yields a modest
          aggregate gain and a dramatic redistribution (B2):}  At the system
          level, REC introduction in B2 raises combined profit by €3{,}244
          (+5.6\,\%).  However, this aggregate improvement masks a large
          intra-market redistribution: Supplier~B (prosumers) gains €10{,}436
          while Supplier~A (consumers) loses €7{,}192.

    \item \textbf{Shared energy volumes are physical-layer invariant:}
          A2 and B2 share near-identical annual shared energy totals
          (75.88\,MWh and 74.40\,MWh).  The physical sharing is determined
          by PV generation and community load profiles, not by supplier structure.
          The financial value of this shared energy, however, is profoundly
          sensitive to how suppliers and balancing groups are arranged.

    \item \textbf{Balancing market effects dominate in two scenarios (A1,
          B2):}  Both A1 and B2 feature large balancing income that outweighs
          structural factors.  In A1, this arises from systematic under-
          forecasting in a predominantly-short system; in B2, it arises from
          REC-induced forecast-actual divergence in Supplier~B's prosumer BG.
          These effects may not persist in future years with different system
          conditions, making them unreliable as long-run profit drivers.

    \item \textbf{The optimal supplier configuration depends on REC strategy:}
          Without a REC, the single-supplier mandate (A1) maximises profit.  With
          a REC, the multi-supplier structure (B2) yields higher aggregate income.
          This inversion suggests that energy policy decisions on supplier
          competition and community energy simultaneously should account for their
          interaction: mandating a single supplier in a high-REC-penetration
          environment could systematically disadvantage suppliers compared to a
          competitive multi-supplier market.
\end{enumerate}

These findings establish the quantitative baseline for Scenario~C
(Battery Optimisation), where a two-stage hierarchical dispatch model will
be evaluated for its ability to further improve both shared energy and supplier
balancing performance beyond the levels achieved in A2 and B2.
